\section*{Limitations}
\label{sec:limitations}

\paragraph{Model coverage.} We evaluate seven API-accessible models from six providers (four proprietary, three open-source; two in replication, Appendix~\ref{app:model_expansion}). Smaller open-source models (7B--70B parameters) may show different vulnerability patterns, and our findings may not generalize to models with fundamentally different training approaches.

\paragraph{Domain scope.} $\tau^2$-bench covers retail and airline customer service. Policy ambiguity effects in other domains (healthcare, finance, legal) remain to be tested, though the within-study cross-domain consistency is encouraging: failure mode distributions share the same top-two modes (assumption-based action and unauthorized escalation) across domains, with individual mode base rates differing by up to 12.5 pp, likely reflecting domain policy structure rather than methodological artifacts.

\paragraph{Synthetic ambiguity.} Our matched-pair design requires synthetic construction of ambiguous variants, which may not capture the full distribution of naturally occurring policy ambiguities. We mitigate this through expert review and a realism validation comparing ambiguous clauses against original $\tau^2$-bench policy text. Ambiguous clauses are indistinguishable from originals in lexical diversity (TTR, KS $p = 0.89$) and word length ($p = 0.09$). Readability is slightly higher (Flesch--Kincaid $d = 0.29$, $p = 0.006$) and sentence length moderately higher ($d = 0.74$), both reflecting the structural demands of encoding ambiguity in standalone clauses. Across matched pairs, disambiguation preserves readability (Flesch--Kincaid $d = 0.14$, $p = 0.10$), confirming that minimal edits change semantic precision without altering linguistic register.

\paragraph{Cross-judge asymmetry.} Although our dual-judge calibration (\S\ref{sec:method}) shows no systematic bias ($p = 0.606$) and human validation confirms the LLM judge operates within inter-annotator disagreement (human--LLM $\kappa = 0.733$ vs.\ inter-annotator $\kappa = 0.798$), each agent model is evaluated by a single cross-judge, so judge identity remains confounded with model identity. Per-type reliability varies: authorization scope has the lowest judge agreement ($\kappa = 0.385$), so the absolute violation rates for this type carry more uncertainty than for incompleteness ($\kappa = 0.875$). Human--LLM agreement shows the same pattern (auth.\ scope $\kappa = 0.644$ vs.\ lexical $\kappa = 0.904$; Appendix~\ref{app:per_type_agreement}). Within-judge-group analyses confirm C1 (both $p < 0.001$) but neither subgroup reaches significance for C3 individually (GPT-4.1--judged: $p = 0.059$; GPT-5.4--judged: $p = 0.157$), leaving open the possibility that the pooled C3 result partly reflects judge-level variance.
Cross-provider test-retest reliability remains unmeasured; ICC $= 0.919$ covers only OpenAI models. Cross-model type ranking concordance is moderate (five-model $W = 0.398$, $p = 0.072$; seven-model $W = 0.353$, $p = 0.038$; Appendix~\ref{app:kendall_w}), with the marginal five-model result reflecting the model$\times$type interaction documented in C4. Within-judge-group type rankings are reported in Appendix~\ref{app:judge_group_ranking}.

\paragraph{Failure mode analysis.} The failure mode taxonomy (\S\ref{sec:analysis}) was coded across all five models (614 traces) using LLM-based coders with substantial inter-coder agreement (Gwet's AC$_1$ $= 0.62$; Cohen's $\kappa = 0.52$). The moderate $\kappa$ reflects the prevalence paradox rather than systematic disagreement, but per-category reliability for minority categories remains limited (e.g., arbitrary rule selection $\kappa = 0.29$). Human validation of the coding scheme would further strengthen these findings.

\paragraph{Incompleteness confound.} As discussed in \S\ref{sec:results}, incompleteness resolution necessarily adds policy content, creating an information asymmetry. Our logistic regression controls and the coreferential lower-bound analysis mitigate this confound, but a future design that orthogonalizes information quantity from ambiguity resolution would strengthen causal claims for this type specifically.

\paragraph{Severity threshold.} Our primary analysis uses moderate-or-above severity. Switching to any-violation (including minor) yields $\Delta = 38.2$ pp (vs.\ 34.0 pp), with all four claims (C1--C4) remaining significant. Type rankings under the two thresholds are concordant (Kendall $\tau = 0.73$), with incompleteness ranking first under both.

\paragraph{Per-type reliability.} Per-type human--LLM $\kappa$ estimates (Appendix~\ref{app:per_type_agreement}) are based on small cell sizes ($n = 32$--$34$), yielding wide confidence intervals (scopal: $[0.00, 1.00]$; coreferential: $[0.00, 0.84]$). These estimates should not be over-interpreted at the individual type level.

\section*{Data and Code Availability}
The 300 clause pairs, evaluation pipeline, analysis scripts, and raw episode traces will be released upon publication at \url{https://github.com/anonymous/policy-ambiguity-safety}.
